% paper/sections/06d_viscous_3layer.tex
%
% §6.3  粘性項：3 層構成と μ∇²u 禁止 (CHK-223: §6 per-term spatial 配下)
%        変粘性二相流で μ∇²u が物理的に誤りである理由と,
%        第 A/B/C 層分離による正しい離散化.
% §6.4  ρ, μ 物性値補間規則 (CHK-223 Phase 6 NEW: ρ/μ prescription co-location)
%        sec:rho_mu_interpolation; ρ=ρ_g+(ρ_l-ρ_g)ψ 代数更新 + 面 μ 算術/調和平均.
%        implicit-BDF2 Helmholtz + デフェクト補正は §7.4 (sec:time_viscous) に集約.
%
% ═══════════════════════════════════════════════════════════════════════════════
\subsection{粘性項：3 層構成と変粘性対応}
\label{sec:viscous_3layer}
% ═══════════════════════════════════════════════════════════════════════════════

\noindent\textbf{連続定式化との関係：}
本節は §~\ref{sec:onefluid}（§2 One-Fluid NS，式~\eqref{eq:NS_onefluid}）の
粘性項 $\bnabla\cdot(2\mu\bm{D})$ を，変粘性二相流の第 A/B/C 層構造として
離散化する設計を示す．界面の応力ジャンプ条件（式~\eqref{eq:jump_stress}：
$[-p\mathbf{I}+\mu_k(\bnabla\bu+\bnabla\bu^T)]_\Gamma\!\cdot\hat{\bm n}_{lg} = \sigma\kappa_{lg}\hat{\bm n}_{lg}$；接線成分は $[\tau_{xy}]_\Gamma=0$）が
3 層分解の設計根拠であり，$\mu\nabla^2\bu$ 直接離散が禁じられる物理的理由となる
（変粘性ジャンプの GFM/IIM 系の取り扱いは~\cite{Kang2000,Sussman2000} を参照）．

本節では変粘性二相流の粘性項を
\textbf{第 A 層（速度勾配生成）・第 B 層（応力テンソル組立）・
第 C 層（応力発散形）}の 3 層に分解し，
各層の独立性を維持することで
$[\tau_{xy}]_\Gamma=0$ を破る離散上の非対称を避ける設計を示す．
これは明示的な界面ジャンプ条件の強制ではなく，
One-Fluid 拡散界面での応力評価をせん断応力連続条件と整合させるための
A3 規律である．

% -------------------------------------------------------------------------------
\subsubsection{\texorpdfstring{なぜ $\mu\nabla^2 u$ が誤りか}{なぜ μ∇²u が誤りか}}
\label{sec:viscous_mu_lap_forbidden}
% -------------------------------------------------------------------------------

非圧縮条件 $\bnabla\cdot\bu=0$ 下で：
\begin{equation}
  \bnabla\cdot(2\mu\mathbf{D})
  = \mu\nabla^2\bu + 2\mathbf{D}\cdot\bnabla\mu,
  \label{eq:viscous_exact_decomposition}
\end{equation}
ここで $\mathbf{D} = \tfrac{1}{2}(\bnabla\bu + \bnabla\bu^\top)$
は歪速度テンソル．
第 2 項 $2\mathbf{D}\cdot\bnabla\mu$ は\textbf{交差微分寄与}で，
$\bnabla\mu\neq 0$ のとき常に非零．
二相流では $\mu$ が界面で跳躍（水-空気で $\mu_l/\mu_g\approx 55$）するため
$\bnabla\mu$ は $\delta$ 関数を含み，交差微分寄与は\textbf{特異で無視不可}．

\noindent\textbf{$\mu\nabla^2\bu$ を使うと}：
\begin{itemize}
  \item $\mu$ の空間変動から生じる接線力を欠落
  \item せん断応力連続 $[\tau_{xy}]_\Gamma = 0$ を満たさない
  \item 界面に沿ったスプリアス接線速度を生成
  \item 高粘性比流でエネルギー散逸が崩壊（$\bu\cdot\mathbf{F}_\mu$ が
        正の値を取りうる）
\end{itemize}

% -------------------------------------------------------------------------------
\subsubsection{3 層構成}
\label{sec:viscous_three_layers}
% -------------------------------------------------------------------------------

\noindent\textbf{第 A 層 — 速度勾配生成}：
コロケート主変数（\S\ref{sec:collocate_layout}）から導出される
\textbf{応力テンソル派生配置}で
4 種の勾配を同一セル中心配置で評価する：
\begin{align}
  (\partial_x u)_{i,j} &\quad\text{セル中心（コロケート $u_{i,j}$ から評価）}, \label{eq:viscous_layerA_1}\\
  (\partial_y v)_{i,j} &\quad\text{セル中心（コロケート $v_{i,j}$ から評価）}, \label{eq:viscous_layerA_2}\\
  (\partial_y u)_{i,j} &\quad\text{セル中心（コロケート $u_{i,j}$ から評価）}, \label{eq:viscous_layerA_3}\\
  (\partial_x v)_{i,j} &\quad\text{セル中心（コロケート $v_{i,j}$ から評価）}. \label{eq:viscous_layerA_4}
\end{align}
第 A 層は領域内部セルで CCD を使用し，界面帯では界面法線方向を
相内 2 次中心 / 片側差分で評価する（\S\ref{sec:viscous_ccd_placement}）．
これは代替スキームへの後退ではなく，粘性ジャンプで生じる
$\partial_n\bu$ の折れを跨がないための応力評価則である．
接線方向が滑らかな場合は CCD を保持できる．

\noindent\textbf{第 B 層 — 応力テンソル組立}：
\begin{align}
  \tau_{xx,i,j} &= 2\mu_{i,j}\,(\partial_x u)_{i,j}, \quad
  \tau_{yy,i,j} = 2\mu_{i,j}\,(\partial_y v)_{i,j}, \label{eq:viscous_layerB_diag}\\
  \tau_{xy,i,j} &= \mu_{i,j}\bigl[(\partial_y u) + (\partial_x v)\bigr]_{i,j}. \label{eq:viscous_layerB_off}
\end{align}
対角成分 $\tau_{xx},\tau_{yy}$ と非対角 $\tau_{xy}=\tau_{yx}$ は
同一セル中心配置で共有し，x-/y-運動量方程式から同じ配列を参照する．

\noindent\textbf{第 C 層 — 応力発散形}：
非一様格子の格子間隔は二系統あり区別する：
\textbf{セル中心間距離} $\Delta x_{i+1/2}\equiv x_{i+1}-x_i$ と
\textbf{セル幅} $\Delta x_i\equiv x_{i+1/2}-x_{i-1/2}$
（$y$ 方向も同様）．第 C 層は応力テンソルに対して
\textbf{一回だけ}物理座標差分を適用する：
\begin{align}
  (\mathcal{F}_\mu)_x\big|_{i,j}
  &= (D_x^{C}\tau_{xx})_{i,j} + (D_y^{C}\tau_{xy})_{i,j}, \label{eq:viscous_layerC_x}\\
  (\mathcal{F}_\mu)_y\big|_{i,j}
  &= (D_x^{C}\tau_{yx})_{i,j} + (D_y^{C}\tau_{yy})_{i,j}. \label{eq:viscous_layerC_y}
\end{align}
$D_\alpha^{C}$ はセル中心応力配列に作用する物理座標差分であり，
界面法線方向では相内 2 次中心 / 片側差分，滑らかな接線方向では
CCD 派生差分を保持できる．
$(\mathcal{F}_\mu)_{i,j}$ は本稿 7 段手順の第5段
（NS 予測子内の implicit-BDF2 粘性 Helmholtz/DC）の高次残差に入る．
第 B/C 層は常に応力発散形；CCD 派生差分を用いるのは
第 A 層の領域内部部分と，界面帯で滑らかと判定された接線方向に限る．
界面法線方向では相内ステンシルを選び，勾配の折れを跨ぐ CCD は使わない．

% -------------------------------------------------------------------------------
\subsubsection{\texorpdfstring{共有せん断応力：$\tau_{xy}$ を第一級要素として扱う}{共有せん断応力：τ\_{xy} を第一級要素として扱う}}
\label{sec:viscous_corner}
% -------------------------------------------------------------------------------

せん断応力 $\tau_{xy}$ はセル中心配置で\textbf{一回だけ定義}し，
x-/y-運動量方程式の両方で同じ配列を参照する．
面/コーナー配置を持つスタガード応力配置へ拡張する場合は，
同じ共有原則の下で $\mu$ をその配置位置へ低次平均する：
\begin{equation}
  \mu_{i+1/2,j+1/2}^\text{arith}
  = \tfrac{1}{4}(\mu_{i,j} + \mu_{i+1,j} + \mu_{i,j+1} + \mu_{i+1,j+1}),
  \label{eq:viscous_mu_corner_arith}
\end{equation}
比較用の調和平均も定義できるが，§\ref{sec:onefluid} の CLS 拡散界面モデルでは
算術平均を基準とする．調和平均は鋭い界面の直列抵抗モデルを仮定する感度評価用であり，
高粘性相の寄与を過小評価し得る：
\begin{equation}
  \mu_{i+1/2,j+1/2}^\text{harm}
  = \biggl[\tfrac{1}{4}\bigl(\mu_{i,j}^{-1} + \mu_{i+1,j}^{-1} + \mu_{i,j+1}^{-1} + \mu_{i+1,j+1}^{-1}\bigr)\biggr]^{-1}.
  \label{eq:viscous_mu_corner_harm}
\end{equation}

標準のコロケート応力配置では，この拡張平均を既定経路には入れず，
セル中心 $\mu_{i,j}$ で構成した共有 $\tau_{xy,i,j}$ を用いる．
同じ式で別々に計算すると数値非対称が入り，離散エネルギー評価が崩れる．

\noindent\textbf{エネルギー評価（理想化）}：
周期境界または境界寄与が消える一様格子で，第 C 層の差分が
第 A 層の勾配の負随伴として働く場合，応力発散形と共有 $\tau_{xy}$ により
\begin{equation}
  \langle\bu,\mathbf{F}_\mu\rangle_h
  = -2\sum_{i,j}\mu_{i,j}\bigl[(\partial_x u)^2+(\partial_y v)^2\bigr]
  - \sum_{i,j}\mu_{i,j}\bigl[(\partial_y u) + (\partial_x v)\bigr]^2_{i,j}
  \le 0.
  \label{eq:viscous_energy_identity}
\end{equation}
対角成分の係数 $-2$ は $\tau_{xx}=2\mu\partial_x u$ の因子 $2$ から，
非対角成分の係数 $-1$ は $\tau_{xy}\cdot(\partial_y u + \partial_x v)
= \mu(\partial_y u + \partial_x v)^2$ から導かれる．
等号は $\mathbf{D}(\bu)\equiv 0$（歪速度テンソル恒等零；
剛体並進・剛体回転で成立する）でのみ成立．
非一様格子，片側境界，および界面帯の法線/接線方向の低次化では
\eqref{eq:viscous_energy_identity} を厳密恒等式とは主張せず，
共有 $\tau_{xy}$ と単一の応力発散を正の人工生成を避けるための設計条件として用いる．

% -------------------------------------------------------------------------------
\subsubsection{CCD 配置：領域内部の第 A 層のみ}
\label{sec:viscous_ccd_placement}
% -------------------------------------------------------------------------------

界面帯では $u\in C^0$（速度連続）だが $\bnabla u\notin C^0$
（粘性跳躍による勾配キンク）．
CCD の $\mathcal{O}(h^6)$ 精度はステンシル支持で場が滑らかであることを前提とするため
キンクを跨ぐ CCD ステンシルは $\partial_n u$ の不連続を増幅し
$\mathcal{O}(1)$ 勾配誤差を生む．

\noindent\textbf{界面帯の定義}：
$|\phi|\le 3h$ を満たすセル / ステンシルは界面帯．
また，ステンシル内に異相のセルを含む場合（相指標の不一致）も界面帯として扱う．

\noindent\textbf{演算子割当}：
\begin{center}
\begin{tabular}{@{}lll@{}}
\hline
領域 & 第 A 層 & 第 B/C 層 \\
\hline
領域内部（界面から遠い）& CCD & 応力発散形 \\
界面帯 & 法線方向は相内 2 次中心 / 片側；接線方向は滑らかなら CCD & 応力発散形 \\
\hline
\end{tabular}
\end{center}

\noindent\textbf{$\phi/\psi$ 判定付き領域内部 Laplacian 縮約}：
領域内部では $\bnabla\mu=0$ のため \eqref{eq:viscous_exact_decomposition} は
$\mu\nabla^2\bu$ に簡約する．
$\psi$ による領域分割で
\begin{equation*}
  L_\mu\bu
  = \mathbb{1}_{\Omega_{\mathrm{I}}}\cdot\mu\Delta_h^{\CCD}\bu
  + \mathbb{1}_{\Omega_\Gamma}\cdot\bnabla\cdot(2\mu\mathbf{D}_h)
\end{equation*}
で領域内部経路を軽量化できる．
$\Omega_\Gamma = \{\psi\in(\varepsilon_\psi, 1-\varepsilon_\psi)\}$
は界面帯指標の一つであり，標準判定は
$|\phi|\le 3h$ と相指標の不一致を優先する．

% ─────────────────────────────────────────────────────────────────────────────
% CHK-223 Phase 3.5: Helmholtz + デフェクト補正は時間積分であるため
% §7.4 (sec:time_viscous) に集約．本ファイルは空間離散化 (第 A/B/C 層
% + 共有 tau_xy + 界面帯ステンシル設計) のみで完結．
% ─────────────────────────────────────────────────────────────────────────────
\medskip
\noindent\textbf{時間積分との結合：}
本節で確立した空間第 A/B/C 層と共有 $\tau_{xy}$ 評価は時間積分のいずれの形式
（陽的 / IMEX / implicit-BDF2）にも入力される．
本稿の implicit-BDF2 では Helmholtz 系 \eqref{eq:helmholtz_implicit_bdf2} の
高次残差 $A_H$ に本節の全応力演算子を用い，
低次補正 $A_L$ は同じ物性値・境界位相・界面帯を共有する
（\S~\ref{sec:time_viscous}）．

% ═══════════════════════════════════════════════════════════════════════════════
% §6.4  ρ, μ 物性値補間規則 (CHK-223 Phase 6: prescription co-location)
% ═══════════════════════════════════════════════════════════════════════════════
\subsection{\texorpdfstring{$\rho, \mu$ 物性値補間規則}{ρ, μ 物性値補間規則}}
\label{sec:rho_mu_interpolation}
% ═══════════════════════════════════════════════════════════════════════════════

\noindent\textbf{設計動機：}
密度 $\rho$ と粘性 $\mu$ は界面で段差状分布（$\rho_g/\mu_g \to \rho_l/\mu_l$）を持つ．
CCD/UCCD6 等の高次差分演算子を直接適用すると Gibbs 振動を生じ，
$[\rho_\text{min},\rho_\text{max}]$ 範囲外の物性値が現れ得る
（\S\ref{sec:scheme_per_variable}，§6.0）．
本節では物性値は常に\textbf{低次代数更新}で評価する規則を確立する．
時間積分章 §7 はこの §6.4 の規則を参照して利用する．

% -------------------------------------------------------------------------------
\subsubsection{\texorpdfstring{セル中心 $\rho, \mu$：CLS $\psi$ からの代数更新}{セル中心 ρ, μ：CLS ψ からの代数更新}}
\label{sec:rho_mu_cell_center}
% -------------------------------------------------------------------------------

CLS 相指標 $\psi^{n+1}\in[0,1]$ から物性値を線形内挿で得る：
\begin{align}
  \rho_{i,j}^{n+1} &= \rho_g + (\rho_l - \rho_g)\,\psi_{i,j}^{n+1},
  \label{eq:rho_update}\\
  \mu_{i,j}^{n+1} &= \mu_g + (\mu_l - \mu_g)\,\psi_{i,j}^{n+1}.
  \label{eq:mu_update}
\end{align}
$\psi$ は CLS 段階 A--F（\S\ref{sec:cls_stage_sequence}，§5.2）で確定し，
段階 F の最終質量閉鎖・値域制限後の値を用いる．
$\rho^{n+1}$ は時間積分 §7 の運動量予測子・PPE・浮力体積力に直接代入される．

\noindent\textbf{線形内挿の正当性：}
$\rho(\psi) = \rho_g + (\rho_l-\rho_g)\psi$ は CLS の連続定式化
（\S\ref{sec:cls_advection}，§3.2）で導出された関係そのもので，
$\psi\in[0,1]$ なら $\rho\in[\rho_g,\rho_l]$ が代数的に保証される．
高次補間（CCD 等）は Gibbs 振動で範囲外値（$\rho<\rho_g$ または $\rho>\rho_l$）を
生じうるため禁止．

% -------------------------------------------------------------------------------
\subsubsection{\texorpdfstring{面上の $\mu$：算術平均 vs 調和平均}{面上の μ：算術平均 vs 調和平均}}
\label{sec:mu_face_interpolation}
% -------------------------------------------------------------------------------

\textbf{面上の $\mu$} は面フラックス型の粘性評価や
implicit-BDF2 Helmholtz の低次補正演算子（§\ref{sec:time_viscous}，§7.4）へ
拡張する場合に必要となる補助量である．
本稿の §6.3 A3 経路はセル中心 $\mu_{i,j}$ で
\eqref{eq:viscous_layerC_x},\eqref{eq:viscous_layerC_y} を組むため，
面上の $\mu$ は既定の応力発散経路ではなく補助規則として位置付ける．
2 種の評価方法：

\noindent\textbf{算術平均}：
\begin{equation}
  \mu_{i+1/2,j}^\text{arith} = \tfrac{1}{2}(\mu_{i,j} + \mu_{i+1,j}).
  \label{eq:mu_face_arith}
\end{equation}

\noindent\textbf{調和平均}：
\begin{equation}
  \mu_{i+1/2,j}^\text{harm} = \frac{2\mu_{i,j}\mu_{i+1,j}}{\mu_{i,j}+\mu_{i+1,j}}.
  \label{eq:mu_face_harm}
\end{equation}

\noindent\textbf{選択指針}：
\begin{itemize}
  \item CLS 拡散界面の基準経路では，粘性は $\mu(\psi)$ の体積平均として扱うため
        算術平均を用いる（§\ref{sec:onefluid}）．
  \item 調和平均は，鋭い界面の直列抵抗則を仮定する感度評価または
        代替配置の比較用に限定する．
  \item 高粘性比の CLS 拡散界面で調和平均を既定にすると，
        高粘性相の寄与を過小評価し得るため，本稿の基準規則にはしない．
\end{itemize}

% -------------------------------------------------------------------------------
\subsubsection{\texorpdfstring{コーナー配置の $\mu$：4 点平均}{コーナー配置の μ：4 点平均}}
\label{sec:mu_corner_interpolation}
% -------------------------------------------------------------------------------

非対角応力 $\tau_{xy}$ を保持する\textbf{コーナー} $(i+1/2, j+1/2)$ では
4 点平均（算術 \eqref{eq:viscous_mu_corner_arith} または調和
\eqref{eq:viscous_mu_corner_harm}）を用いる．
これはスタガード/コーナー応力配置へ拡張する場合の規則であり，
標準のコロケート応力発散ではセル中心 $\mu$ と共有 $\tau_{xy,i,j}$ を用いる．
共有配列原則は \S\ref{sec:viscous_corner}（§6.3）を参照．

% -------------------------------------------------------------------------------
\subsubsection{時間積分との結合順序}
\label{sec:rho_mu_time_integration}
% -------------------------------------------------------------------------------

$\rho^{n+1}, \mu^{n+1}$ は時間ステップの処理順序の中で
\textbf{CLS 段階 A--F 後・運動量予測子前}に確定する：
\begin{enumerate}
  \item 段階 F の最終質量閉鎖後の $\psi^{n+1}$ から \eqref{eq:rho_update}, \eqref{eq:mu_update}
        でセル中心 $\rho^{n+1}, \mu^{n+1}$ を更新．
  \item 本稿のセル中心応力発散ではセル中心 $\mu^{n+1}$ を
        第 B 層の応力組立に直接用いる．
  \item 面/コーナー配置を持つ拡張配置では，
        \eqref{eq:mu_face_arith}--\eqref{eq:mu_face_harm} または
        \eqref{eq:viscous_mu_corner_arith}--\eqref{eq:viscous_mu_corner_harm}
        で低次平均を評価する．
\end{enumerate}
詳細な速度--PPE 結合順序は
\S\ref{sec:cls_velocity_consistency_v7}（§7.3）を参照．

% -------------------------------------------------------------------------------
\subsubsection{失敗モードと回避}
\label{sec:rho_mu_failure_modes}
% -------------------------------------------------------------------------------

\begin{itemize}
  \item \textbf{$\rho, \mu$ に CCD を直接適用} → Gibbs 振動 → 範囲外値．
        必ず代数更新 \eqref{eq:rho_update}, \eqref{eq:mu_update}．
  \item \textbf{CLS 拡散界面で面上の $\mu$ の調和平均を既定化} → 高粘性相の寄与を過小評価
        → 粘性散逸の過小評価．
        本稿の基準経路では算術平均，調和平均は感度評価に限定．
  \item \textbf{$\tau_{xy}$ を x-/y- 運動量で別々に計算} → 数値非対称
        → エネルギー評価と安定性が破綻．
        共有配列原則（\S\ref{sec:viscous_corner}）．
  \item \textbf{$\psi$ 値域制限前に物性値更新} → $\rho<\rho_g$ または $\rho>\rho_l$．
        段階 B/F の閉鎖と $\psi\in[0,1]$ 値域制限後に更新．
\end{itemize}

% ═══════════════════════════════════════════════════════════════════════════════
